% GENERATED FILE. Edit canonical lessons, then run npm run generate:books.
% canonical-source-hash: fnv1a64:6bc54366f7b2e32b
% canonical-lessons: ES-C272-es-que, ES-C272-si-claro, ES-C272-tomar-el-pelo, ES-C272-mira, ES-C272-practice

\chapter{Reading Between the Lines --- Imply, Invert, Look It Up, Re-tune}
\label{ch:entrelineas}
\hlchaptermodality{272}

\begin{chapteropening}
\textbf{By the end of this chapter:} \emph{I can tell the four ways a Spanish sentence means more than it says, and say which strategy each one asks of the listener.}

Complete ES-C272-practice: name the listener strategy each of the four devices requires, and which one a written text hides.
\end{chapteropening}

\section[es que…]{es que… — the refusal that never says no}

\label{lesson:ES-C272-es-que}



\begin{quote}
Everything so far has meant what it said. This chapter is about Spanish that means MORE than it says, and this is the commonest case.
\end{quote}

\subsection*{What to know first}
\begin{itemize}
\raggedright
  \item supongo que — the hedge that claims an inference.
\end{itemize}

\begin{sounds}
\begin{itemize}
\raggedright
  \item \emph{es ke}, run together, and usually trailing off rather than finishing. $\to$ reference
\end{itemize}
\end{sounds}

\begin{cousinweb}
\textbf{es que…} — literally \textquotedblleft{}it is that…\textquotedblright{}, and functionally almost never that.

Asked to do something you cannot do, a Spanish speaker very often begins \emph{es que…} and then gives a circumstance rather than an answer. \emph{¿Vienes? — Es que mañana trabajo.} Nothing in that sentence is a refusal. Every listener hears one.

The words supply a reason; the \textbf{refusal is left for you to work out}. That gap between what is said and what is meant has a name in linguistics — an \textbf{implicature} — and Spanish leans on it harder than English does.
\end{cousinweb}

\begin{grammarlens}[title={What you've built}]
This is the first thing in the book you cannot translate word by word and still understand. \emph{It is that tomorrow I work} is not English, and even repaired to \emph{it's just that I work tomorrow} it loses the refusal.

Two consequences worth carrying. When you hear \emph{es que}, expect the sentence after it to be a reason for a \textquotedblleft{}no\textquotedblright{} that will not be spoken. And when you need to decline something yourself, this is the polite instrument for it.
\end{grammarlens}

\subsection*{Your turn}
Say these aloud:

\begin{itemize}
\raggedright
  \item \emph{es que…} and finish it with a reason
  \item what a listener would understand from it, in one word (\emph{no})
  \item the bare refusal instead, and notice how much blunter it lands
\end{itemize}

\begin{tcolorbox}[breakable,colback=teal!4,colframe=teal!35!black,title={Before you move on}]
What is the gap between what a sentence says and what it means called? (An implicature.) What does \emph{es que} usually introduce? (A reason standing in for a refusal.)
\end{tcolorbox}
\section[sí, claro]{sí, claro — when the words say yes and the speaker means no}

\label{lesson:ES-C272-si-claro}



\begin{quote}
Harder than an implicature: a sentence whose meaning is the OPPOSITE of its words.
\end{quote}

\subsection*{What to know first}
\begin{itemize}
\raggedright
  \item es que… — meaning more than you say.
\end{itemize}

\begin{sounds}
\begin{itemize}
\raggedright
  \item Everything here is in the delivery. Said briskly and warmly, \emph{sí, claro} is plain agreement. Said slowly, flatly, with the \emph{claro} dragged out, it is the opposite. $\to$ reference
\end{itemize}
\end{sounds}

\begin{cousinweb}
\textbf{sí, claro} — \textquotedblleft{}yes, of course\textquotedblright{} — or \textquotedblleft{}yeah, right\textquotedblright{}, depending entirely on how it is said.

\emph{Claro} is Latin \textbf{clarus}, \textquotedblleft{}bright, clear\textquotedblright{} — the same word as English \textbf{clear} and the name \textbf{Clara}. To say \emph{claro} is to say \emph{that is obvious}.

Which is exactly what makes it available for irony. A word meaning \textquotedblleft{}obviously\textquotedblright{} is the perfect thing to say about something you find not obvious at all, and the gap between the word and the tone is where the sarcasm lives.
\end{cousinweb}

\begin{grammarlens}[title={What you've built}]
Irony is not a Spanish trick; English does the same with \emph{sure} and \emph{of course}. What is worth knowing is that \textbf{the marker is carried in the voice, not in the words} — so a written \emph{sí, claro} is genuinely ambiguous, and Spanish readers resolve it from context exactly as you do in English.

If you are not certain which one you are hearing, you are not failing at Spanish. The sentence is ambiguous by design.
\end{grammarlens}

\subsection*{Your turn}
Say these aloud:

\begin{itemize}
\raggedright
  \item \emph{sí, claro} meaning it — brisk and warm
  \item \emph{sí, claro} not meaning it — slow and flat
  \item which of the two a written text cannot tell you
\end{itemize}

\begin{tcolorbox}[breakable,colback=teal!4,colframe=teal!35!black,title={Before you move on}]
What does \emph{claro} mean literally, and which English word is it? (Bright, clear — \emph{clear}.) Where is the irony carried? (In the voice, not the words.)
\end{tcolorbox}
\section[tomar el pelo]{tomar el pelo — four words you know, meaning something else entirely}

\label{lesson:ES-C272-tomar-el-pelo}



\begin{quote}
You know every word in this phrase. That will not help you.
\end{quote}

\subsection*{What to know first}
\begin{itemize}
\raggedright
  \item tomar — \textquotedblleft{}to take\textquotedblright{}, which you already own.
  \item sí, claro — the irony marker just before this.
\end{itemize}

\begin{sounds}
\begin{itemize}
\raggedright
  \item \emph{to-MAR el PE-lo}. $\to$ reference
\end{itemize}
\end{sounds}

\begin{cousinweb}
\textbf{tomar el pelo} — \textquotedblleft{}to tease someone, to have someone on.\textquotedblright{}

Word by word it is \emph{to take the hair}, and word by word it tells you nothing. \emph{Pelo} is Latin \textbf{pilus}, \textquotedblleft{}a hair\textquotedblright{} — the same root inside English \textbf{pile} (the raised surface of a fabric) and \textbf{depilate}.

Where the phrase comes from is genuinely uncertain, and the tidy stories you will find online are folk etymology. What is certain is that its meaning cannot be reached from its parts, and that is the definition of an \textbf{idiom}.

English does the same to Spanish learners with \emph{pull someone's leg}, which has no better-attested origin.
\end{cousinweb}

\begin{grammarlens}[title={What you've built}]
An idiom is the limit case of this chapter. An implicature asks you to infer; irony asks you to invert; an idiom asks you to \textbf{look it up}, because there is nothing in the words to work with.

That is worth stating plainly rather than leaving you to discover it: some Spanish cannot be decoded, only learned. Recognising which is which is the skill this chapter is really teaching.
\end{grammarlens}

\subsection*{Your turn}
Say these aloud:

\begin{itemize}
\raggedright
  \item \emph{tomar el pelo}
  \item what it says word by word (\emph{to take the hair})
  \item whether the origin story you may have heard is reliable (it is not)
\end{itemize}

\begin{tcolorbox}[breakable,colback=teal!4,colframe=teal!35!black,title={Before you move on}]
What does \emph{tomar el pelo} mean? (To tease, to have someone on.) Can its meaning be worked out from its words? (No — that is what makes it an idiom.)
\end{tcolorbox}
\section[mira]{mira — a word that changes the temperature of what follows}

\label{lesson:ES-C272-mira}



\begin{quote}
The last of the four, and the smallest: one word that tells the listener how to take the next sentence.
\end{quote}

\subsection*{What to know first}
\begin{itemize}
\raggedright
  \item tomar el pelo — the idiom before this.
\end{itemize}

\begin{sounds}
\begin{itemize}
\raggedright
  \item \emph{MI-ra}, stress the first syllable. $\to$ reference
\end{itemize}
\end{sounds}

\begin{cousinweb}
\textbf{mira} — \textquotedblleft{}look.\textquotedblright{}

It is the command form of \emph{mirar}, \textquotedblleft{}to look\textquotedblright{} — Latin \textbf{mirari}, \textquotedblleft{}to wonder at, to marry at\textquotedblright{}, which also gave English \textbf{miracle}, \textbf{admire} and \textbf{mirror}.

But dropped at the start of a sentence it has stopped being a command. Nobody is being asked to look at anything. \emph{Mira, no es así} — \textquotedblleft{}look, it isn't like that\textquotedblright{} — uses it to mark what follows as a correction, an intimacy, or an impatience, depending on the voice.

Spanish is full of these: \emph{oye}, \emph{venga}, \emph{bueno}, \emph{hombre}. Every one of them started as a word with a meaning and became a \textbf{signal about the sentence around it}.
\end{cousinweb}

\begin{grammarlens}[title={What you've built}]
This is a \textbf{discourse marker}, and it does not translate — it \emph{positions}. The literal meaning is spent; what is left is a change in temperature.

That completes the chapter's four ways a Spanish sentence can mean more than it says: it can imply (\emph{es que}), it can invert (\emph{sí, claro}), it can be opaque (\emph{tomar el pelo}), or it can be re-tuned by a word that carries no content at all (\emph{mira}).
\end{grammarlens}

\subsection*{Your turn}
Say these aloud:

\begin{itemize}
\raggedright
  \item \emph{mira} as a real command — look at something
  \item \emph{mira} as a marker — \emph{mira, no es así}
  \item two more markers Spanish uses this way
\end{itemize}

\begin{tcolorbox}[breakable,colback=teal!4,colframe=teal!35!black,title={Before you move on}]
What is \emph{mira} literally? (The command \textquotedblleft{}look\textquotedblright{}, from \emph{mirar}.) What does a discourse marker do instead of meaning something? (It positions the sentence — it changes how you are meant to take it.)
\end{tcolorbox}
\section[Practice]{Four ways to mean more than you say}

\label{lesson:ES-C272-practice}



\begin{quote}
Four devices, and the useful skill is telling them apart fast enough to keep up with a speaker.
\end{quote}

\subsection*{Your turn}
For each of the four, say which strategy the listener has to use.

Say these aloud:

\begin{itemize}
\raggedright
  \item \emph{es que mañana trabajo} — and then say what the listener must do (\textbf{infer} — supply the refusal that was not spoken)
  \item \emph{sí, claro}, flat and slow — and what the listener must do (\textbf{invert} — take the opposite of the words)
  \item \emph{me estás tomando el pelo} — and what the listener must do (\textbf{know it already} — nothing in the words will help)
  \item \emph{mira, no es así} — and what the listener must do (\textbf{re-tune} — the marker says how to take what follows, not what it means)
\end{itemize}

Say these aloud:

\begin{itemize}
\raggedright
  \item which of the four a written text hides from you, and why (\emph{sí, claro} — the irony is in the voice)
  \item which of the four you could not have guessed without being taught
\end{itemize}

\begin{tcolorbox}[breakable,colback=teal!4,colframe=teal!35!black,title={Before you move on}]
Name the four strategies a listener uses: infer, invert, look it up, re-tune. Which one is carried in the voice rather than the words? (Irony.) Which one cannot be worked out at all? (The idiom.)
\end{tcolorbox}
